\section{The marked block integral on the original activity addresses}
\label{sec:marked-block-integral}
We now obtain the exact tilted identity used in
Section~\ref{sec:leaf-expansion}. The insertion is transported through
the same integral and the same connected-component decomposition as
the background. There is no division by a forced activity, which may
vanish, and the distinguished activity address remains fixed.

\subsection{The fine inserted density}
Write the fine density as
\[
 \rho(U)=e^{E(U)}\sum_{A\ {\rm admissible}}\prod_{i\in A}K_i(U).
\]
Choose a bounded real measurable gauge-invariant function $h$, with
$|h|\le w$, supported on the eligibility condition of one fixed
activity address. Its dependence on the configuration is confined to
that address's full inspected support. There are two possibilities.
For a regular address $b$, its background regular activity is
identically zero and $h$ vanishes outside the healthy set for its
support. For a forced address $b$, $h$ vanishes whenever that forced
address fails its original eligibility test. Define
\begin{equation}\label{eq:marked-fine-density}
 f_h(U)=e^{E(U)}\sum_{A\ {\rm admissible}} B_h(A,U),\qquad
 B_h(A,U)=
 \begin{cases}
 h(U)\prod_{i\in A}K_i(U),&b\text{ regular},\\
 \mathbf1_{\{b\in A\}}h(U)\prod_{i\in A\setminus\{b\}}K_i(U),
                                      &b\text{ forced}.
 \end{cases}
\end{equation}
The admissibility tests, including exact flag coverage, are those
of the background. In the forced case the address is replaced even
when its original value $K_b(U)$ is zero.

Here and below use the same funded sector data and all-species
integrated product bounds as in
Section~\ref{sec:gaussian-reassembly}. On the fixed finite lattice,
choose summable uniform prices $p_i\ge|K_i(U)|$. In the regular case
the absolute family sum in \eqref{eq:marked-fine-density} is bounded by
$w\prod_i(1+p_i)$. In the forced case replace $p_b$ by $w$ and use
the same product bound. Thus the series converges uniformly and
absolutely. The fine exponent has a bounded measurable envelope, so
$f_h$ is bounded and measurable. Gauge invariance follows termwise
from the invariance of $h$, the activities and the flag tests.
In particular the fine Haar integral and every fixed transition
slice integral exist. This argument does not require a sign for $h$.

\subsection{One marked factor in the Gaussian expansion}
Denote the elementary alphabet of the unmarked expansion by
$\mathcal I$ and the letter assigned to $b$ by $\star$. Its support
is the same enlarged inspected support as the original address.
For one sector, one tag assignment and one old Haar configuration,
let $v_i(z)$ be the elementary factors of the Gaussian expansion.
The marked factor is
\[
 H(z)=h(U(z))\,\zeta_\star(z),\qquad 0\le\zeta_\star\le1,
\]
where $\zeta_\star$ is the larger local fluctuation cutoff on the
marked support. For a finite set of letters $P$ put
\begin{equation}\label{eq:marked-letter-product}
 \Psi_h(P,z)=\mathbf1_{\{\star\in P\}}H(z)
                          \prod_{i\in P\setminus\{\star\}}v_i(z).
\end{equation}
The cutoff makes the marked factor part of the same local integral
construction. On the nonzero support of the genuine common chart
cutoff it equals one: the stronger common small-fluctuation bound
places each marked coordinate in its plateau. If the common cutoff
is zero, both sides of the chart identity are zero. Hence inserting
$\zeta_\star$ changes none of the represented fine integral.

The regular and forced addresses give two slightly different
expansion identities. At a blank regular address the old
exponential-minus-one factor $v_\star$ is zero. Every old term
containing that letter therefore vanishes. Insert $\star$ into each
remaining finite letter set and replace its zero factor by $H$.
Erasing $\star$ is the inverse map, so the resulting sum is $H$
times the original sum. In the forced case, a family not containing
the address contributes zero. For a family containing it, erase
just that forced factor and put $H$ in its place. The optional
logarithmic and warning expansions are unchanged. These two
identities give precisely the two cases of
\eqref{eq:marked-fine-density}; there is no extra multiplicity.

To justify integration of these expansions, let $r_i\ge0$ be the
summable elementary integrated prices already established for the
unmarked product construction. Its bound applies to every finite
set of letters, in particular to $P\setminus\{\star\}$. Since
$|H|\le w$, it gives
\begin{equation}\label{eq:marked-integrated-price}
 \int|\Psi_h(P,z)|\,d\nu(z)
 \le \mathbf1_{\{\star\in P\}}w
                            \prod_{i\in P\setminus\{\star\}}r_i,
\end{equation}
where $\nu$ includes the old probability Haar variables and the
Gaussian variables. The right side is summable over $P$: erasure
maps the sets containing $\star$ injectively to all finite sets,
whose product-price sum is at most $\exp(\sum_i r_i)$.
The separate uniform envelope prices give measurability and an
integrable envelope exactly as in the unmarked construction.
Thus Fubini and absolute summation justify every exchange here.

For clarity, the chart and Gaussian normalization in this argument
are unchanged. First split the free Haar coordinates into the
active chart coordinates and their complement. On a nonzero sector
the chart-support lemma puts every active bond in the original
exponential chart. The product change of variables contributes
$c_H^{N_{\rm active}}\prod J(u_i)$; the affine rescaling contributes
$g^{qN_{\rm active}}$; the unused Gaussian coordinates integrate to
one with their normalized density. The quadratic completion and
logarithmic cancellation from
Section~\ref{sec:gaussian-reassembly} then apply with the extra
factor $h(U)$ untouched. Its local cutoff has just been restored
to one on this support. Summing the sector partition and averaging
the original tags consequently recovers $\mathcal T f_h$, with the
same scalar overhead $\mathscr O$ as for $\mathcal T\rho$.

\subsection{The distinguished connected component}
Split a finite elementary letter set into its connected components,
using their full assigned supports. If it contains $\star$, exactly
one component $C$ contains that letter. Formula
\eqref{eq:marked-letter-product} factors into the marked product on
$C$ and the ordinary products on the other components. The Gaussian
cross covariances between their read sets vanish; the Haar and tag
read sets are disjoint. Thus the product integral factors. This is
the same finite-range independence argument used in the unmarked
case, now with one bounded local factor substituted.

Let $a_C^F(V)$ be an ordinary integrated component value and
$b_C^F(V)$ the marked value, set to zero unless $\star\in C$.
The component and flattening maps are inverse on the permitted
families. The sector and eligibility masks are identical on both
sides: if a component fails its local test, its corresponding
ordinary or marked value is zero; on a nonzero family the union of
its inventories is exactly the set of flags. Absolute convergence
from \eqref{eq:marked-integrated-price} therefore proves
\begin{equation}\label{eq:marked-forest-integral}
 \mathcal T f_h(V)=\mathscr O
 \sum_{\mathcal A\ {\rm admissible}}
    \sum_{C\in\mathcal A}b_C^F(V)
                       \prod_{D\in\mathcal A\setminus\{C\}}a_D^F(V).
\end{equation}

The marked value has the same locality in the forest as the ordinary
one. Indeed $h$ reads only its inspected support, and its added
cutoff reads only the corresponding buffered coordinates. Forests
agreeing there give equal values after pulling their free Haar
variables back to the common full Haar space. Their principal
Gaussian covariances agree, so integration gives equality too.
We may therefore average independent forest choices at coarse
sites as in Section~\ref{sec:forest-average-log}. For a compatible
family the marked component and all the ordinary components use
disjoint sets of site choices, so the average of their product is
the product of their averages. Set
\[
             a_C=\mathbb E_F a_C^F,\qquad b_C=\mathbb E_F b_C^F.
\]
Equation~\eqref{eq:marked-forest-integral} holds with these averaged
values. The same absolute majorants apply to this finite average.

Two distinct compatible components cannot both have nonzero marked
value, since both would contain $\star$. Hence, for every real
$\theta$ and every compatible family,
\[
 \prod_{C\in\mathcal A}(a_C+\theta b_C)
 =\prod_{C\in\mathcal A}a_C
  +\theta\sum_{C\in\mathcal A}b_C
                           \prod_{D\ne C}a_D.
\]
The terms of degree at least two vanish individually. Together
with the unmarked identity this proves the exact integral formula
\begin{equation}\label{eq:marked-affine-integral}
 \mathcal T\rho+\theta\mathcal T f_h
   =\mathscr O\sum_{\mathcal A\ {\rm admissible}}
                               \prod_{C\in\mathcal A}(a_C+\theta b_C).
\end{equation}
The right side is affine despite its product notation. The masks,
support assignments, inventories and distinguished fine address
have never changed.

\subsection{Passing the mark through the two regroupings}
For $\theta$ in the interval allowed by the purchased marked prices,
apply the regular cluster expansion and forced exclusion regrouping
to the values $a_C+\theta b_C$. Denote their regular exponential
by $a(\theta)$ and their forced gas by $G_\theta$. The exact
identities of Sections~\ref{sec:forest-average-log}
and~\ref{sec:forced-exclusion} give
\[
 \mathcal T\rho+\theta\mathcal T f_h
                         =\mathscr O a(\theta)G_\theta.
\]
The absolute product and cluster majorants apply uniformly on this
interval. Indeed if $|a_C|\le A_C$, $|b_C|\le D B_C$ and
$|\theta|D\le1$, the affine factor has the joint value/derivative
bound $A_C+2B_C$ of Section~\ref{sec:leaf-expansion}. Products,
cluster coefficients and exponential multiplicities retain that
joint bound by the product rule proved there. This uses the
marked-price purchases; the unmarked prices alone do not supply
the derivative bound.

Finally let $R_0,S_0$ be the two fixed local logarithmic corrections
absorbed into the forced gas in Section~\ref{sec:soft-log-fusion}.
That regrouping uses the same fixed background and gives
$b(\theta)=G_\theta e^{R_0+S_0}$.

We identify the coarse exponent on the same configuration, background,
path and clipping threshold as the integral above. The regular
families have the summability, symmetry and local reset hypotheses of
Section~\ref{sec:physical-reset-reorganization}, including its side-size
and clipping conditions. Let $F$ be the exponent
of the intermediate representation with regular terms frozen, and
let $P$ be the path contribution subtracted with that quadratic
coefficient included in it. We shall prove
\begin{align}
 E_{\rm coarse}+S_0&=F-P,\label{eq:marked-output-exponent-input}\\
 e^{F+R_0-P}&=\mathscr O a(0).\label{eq:marked-regular-accounting-input}
\end{align}
Here are the two calculations. All sums below are over the same
finite coarse lattice and, for the local entries, over absolutely
summable families. Write the frozen exponent as
$F=\tau N-K_\beta-A$, where $N$ is the number of coarse sites,
$K_\beta$ is the retained path exponent before the coupling reset,
and $A$ is the frozen regular sum. In the second regrouping let $v$
be the vacuum subtraction, $d$ the sum of the discarded regular
terms, $a$ the sum of the untagged terms, $H$ the quadratic proxy
sum, and $L_0$ the sum of the path loans. Let $\Delta$ be the
extracted quadratic coefficient. The splitting of the regular
entries and the path ledger take the form
\[
 A=v+\Delta H-d+a,\qquad
 K_{\beta+\gamma+\Delta}-K_\beta+J=P+\Delta H+L_0,
\]
where $J$ is the sum of the lower and upper path corrections added
to the regular terms, and $P$ is $\gamma$ times the masked path
receipt. These are precisely \eqref{eq:reset-first-split} and
\eqref{eq:reset-path-identity}: in that section's notation the present
$(v,H,d,a,J,L_0,\gamma)$ are $(c_0,h,D,R_0,F,J,\alpha)$.
The kinetic difference is $(\gamma+\Delta)\mathcal K$, since only
its coupling changes; the loop term is unchanged. Thus these identities
hold on every field, including unhealthy ones.
The reset uses vacuum $\tau-v/N$, regular sum $a+J$
and logarithmic sum $S_0=d+L_0$. Collecting the terms gives
\[
 E_{\rm coarse}+S_0
 =\tau N-K_\beta-(v+a+\Delta H-d)-P=F-P.
\]
The final regrouping preserves this exponent because it preserves
the sum of the regular entries, leaving the vacuum and path terms
unchanged.

For completeness, the first transfer can be checked directly on its
entries. Let $A_*$ be the original regular sum and $A_{\rm rem}$
the sum after removing its selected finite set $I_0$ of entries.
Distribute each selected masked entry $f_i$ over sites with its
allocation weights $\lambda_{oi}$, for which $\sum_o\lambda_{oi}=1$,
and put $e_o=\sum_{i\in I_0}\lambda_{oi}f_i$. Hence
$A_*=A_{\rm rem}+\sum_o e_o$. Let $c_*$ be the common flat value
of the allocated local reference (independent of $o$ by translation),
let $g_o$ be its local replacement, and write $\chi_o$ for that
replacement's healthy mask. The discrepancy used by the transfer is
\[
 t_o=e_o-c_*-\chi_o g_o.
\]
Thus $A_*=A_{\rm rem}+Nc_*+\sum_o\chi_o g_o+\sum_o t_o$.
This uses the actual masked entries and is valid on unhealthy fields
too. The local replacement makes $t_o$ vanish on its larger healthy
support, as required for absorbing $-t_o$ into the gas.

Write $\langle\cdot\rangle$ for the uniform average over signed
translations of the coarse lattice. The original and remaining
regular sums are invariant under this average: their label
permutations preserve values and masks, and their series converge
absolutely. The new regular sum, vacuum and logarithmic correction are
\[
 A=A_{\rm rem}+\left\langle\sum_o\chi_o g_o\right\rangle,
 \qquad \tau=\tau_*-c_*,
 \qquad R_0=-\left\langle\sum_o t_o\right\rangle.
\]
The averaging factor is the reciprocal of the number of signed
translations, so the constant $Nc_*$ is unchanged. Substitution gives
$F+R_0=N\tau_*-K_*-A_*$. Enumerating the two disjoint label
families changes no sum, so this also proves the identity for the
actual enumerated arrays.

We next identify $A_*$. Write $c_\Gamma(U)$ for the absolutely
summable flag-free cluster logarithms of
Section~\ref{sec:forest-average-log}. Their sum is $\log a(0)$
and $a(0)=\exp(\sum_\Gamma c_\Gamma)>0$.
The first part of the regular alphabet has values $-c_\Gamma$,
transported to fixed addresses with the same full supports.
Its assigned healthy mask changes no value, because $c_\Gamma$
already vanishes off that healthy set. The transport is a bijection,
so this part sums to $-\log a(0)$ without a multiplicity.

The other part consists of the low-path receipts with coefficient
$-\gamma$, the cubic and quartic atoms with coefficient $-\beta$,
and the loop atoms with coefficient $-1$. Each retains its original
mask and signed-channel weight. The low channels sum to $-P$ and
the other three families sum to $-B_*$, where
\[
 B_*=\beta(\mathcal B_3+\mathcal B_4)+\mathcal B_{\rm loop},
 \qquad K_*=K_{\rm base}+B_*.
\]
Each $\mathcal B$ is the signed average of its translated masked
path atoms. Partitioning the path alphabet proves these equalities;
the path bounds give absolute summability. Translation and signed
averaging therefore commute with the sums. The cluster sum is
invariant under the same symmetries by their bijections of cluster
addresses. The further uniform signed-translation average of the
whole regular alphabet leaves its sum unchanged. This proves
$A_*=-\log a(0)-P-B_*$.

Finally let $T_*>0$ be the field-independent factor in
\eqref{eq:reassembly-overhead}, evaluated at the identity with the
same forest and incoming vacuum. The extracted field-dependent part
of the inherited exponent is $\exp(-K_{\rm base})$, whence
$\mathscr O=T_*\exp(-K_{\rm base})$. The vacuum is defined by
$\tau_*=N^{-1}\log T_*$; $N>0$ and $T_*>0$ justify this definition.
Writing $v_*=N\tau_*$ now gives
\[
 F+R_0-P=v_*-K_{\rm base}-B_*-A_*-P
        =v_*-K_{\rm base}+\log a(0).
\]
Exponentiating proves \eqref{eq:marked-regular-accounting-input}.
Every reindexing above preserves the original threshold and path;
in particular the receipt $P$ used in the reset is the one just
summed in the regular alphabet.

Indeed, adding $R_0$ to \eqref{eq:marked-output-exponent-input}
and exponentiating, then using \eqref{eq:marked-regular-accounting-input},
gives
\begin{equation}\label{eq:marked-coarse-normalization}
 e^{E_{\rm coarse}+R_0+S_0}=\mathscr O a(0).
\end{equation}
With this normalization, substitution proves
\[
 \mathcal T\rho+\theta\mathcal T f_h
 =e^{E_{\rm coarse}}\frac{a(\theta)}{a(0)}b(\theta),
\]
which is \eqref{eq:leaf-tilted-identity} for the actual insertion
\eqref{eq:marked-fine-density}. Differentiation is justified by
the joint summable bounds, not by a formal manipulation of the
series. This establishes the exact marked integral underlying that
argument. The remaining quantitative task is to purchase its marked
local derivative prices and principal-term estimate on the selected
history; the density identity alone does not provide those estimates.
